%! Carga de paquetes
\usepackage{amssymb,amsthm,amsmath,mathtools,dsfont}
\usepackage{graphicx,multicol,tikz,accents,multimedia,tcolorbox}
\usepackage[spanish]{babel}
\usepackage[utf8]{inputenc}

%! soluciona potencial problema de fuentes en math mode
%\usefonttheme[onlymath]{serif}
\usefonttheme{professionalfonts} 

%! Personalización de colores
%* Colores institucionales según documento oficial
\definecolor{BlueUSM}{HTML}{004B85}
\newcommand{\BlueUSM}[1]{{\color{BlueUSM}{#1}}}
\definecolor{RedUSM}{HTML}{D60019}
\newcommand{\RedUSM}[1]{{\color{RedUSM}{#1}}}
\definecolor{GreenUSM}{HTML}{008452}
\newcommand{\GreenUSM}[1]{{\color{GreenUSM}{#1}}}
\definecolor{YellowUSM}{HTML}{F7AE00}
\newcommand{\YellowUSM}[1]{{\color{YellowUSM}{#1}}}
\definecolor{BlackUSM}{HTML}{000000}
\newcommand{\BlackUSM}[1]{{\color{BlackUSM}{#1}}}
\definecolor{LightBlackUSM}{cmyk}{0,0,0,0.7}
\newcommand{\LightBlackUSM}[1]{{\color{LightBlackUSM}{#1}}}

%! Beamer configuración básica
\mode<article>{\usepackage{pgf,fancyhdr}}
\mode<presentation> {
\usetheme{JuanLesPins}
% Color del tema. (default es el azul tipo "puc")
\usecolortheme[named=BlueUSM]{structure}
%\setbeamercovered{transparent} % Transparencia.
}
%\setbeameroption{hide notes}
%\setbeameroption{show notes}
% left, bottom, or top. (default is right)
%\setbeameroption{show notes on second screen=⟨location⟩}

%! Traducción de ambientes en Beamer
\newtranslation[to=spanish]{Section}{Secci\'on}
\newtranslation[to=spanish]{Subsection}{Subsecci\'on}